\subsection{Unitary dual of a locally compact group}

Let G be a locally compact topological group. One equips G with a left Haar measure \(\mu\) . For \(p \in [1, +\infty]\) , we denote by \(\mathcal{L}^{p}(\mathrm{G})\) (resp. \(\mathrm{L}^{p}(\mathrm{G})\) ) the space \(\mathcal{L}_{\mathbf{C}}^{p}(\mathrm{G}, \mu)\) (resp. the space \(\mathrm{L}_{\mathbf{C}}^{p}(\mathrm{G}, \mu)\) ). We identify the space \(\mathcal{C}_{b}(\mathrm{G})\) with its image in \(\mathrm{L}^{\infty}(\mathrm{G})\) .

Denote by \(\Delta\) the modular function of G. Recall that \(L^1(G)\) is an involutive Banach algebra whose involution is induced, by passage to quotients, by the map \(f \mapsto f^*\) where \(f^*(y) = \Delta^{-1}(y)\overline{f(y^{-1})}\) for all \(f \in \mathcal{L}^1(G)\) and \(y \in G\) (cf. I, p. 99, example 4). The Banach algebra \(L^1(G)\) admits an approximate unit according to INT, VIII, p. 172, §4, no. 7, prop. 20.

Let \(\varrho\) be a unitary representation of G in a complex Hilbert space E. The map \(f\mapsto\varrho(f)\) is a continuous morphism of involutive algebras from L \(^{1}\) (G) to \(\mathcal{L}\) (E) (lemma 1 of V, p. 401); it is a non-degenerate representation of L \(^{1}\) (G) in E (INT, VIII, p. 139, § 2, no. 7, prop. 10, (i)).

PROPOSITION 12. --- Let E be a complex Hilbert space and \(\widetilde{\pi}\) a morphism of involutive algebras from L \(^{1}\) (G) to \(\mathcal{L}\) (E) . If the representation \(\widetilde{\pi}\) is nondegenerate, then there exists a unique unitary representation \(\pi\) of G in E such that \(\widetilde{\pi}(f)=\pi(f)\) for all \(f\in\mathrm{L}^{1}(\mathrm{G})\) .

Uniqueness follows from INT, VIII, p. 139, § 2, n° 7, cor. 3 of lemma 4.

Let us prove the existence of \(\pi\) . We have \(\| \widetilde{\pi}\| \leqslant 1\) by prop. 2 of I, p. 104. Let F be the subspace of E generated by the vectors of the form \(\widetilde{\pi}(f)x\) for \(f\) in \(\mathrm{L}^1 (\mathrm{G})\) and \(x\) in E; it is dense in E since the representation \(\widetilde{\pi}\) is nondegenerate.

For every compact neighborhood V of e, let \(\varphi_{V}\) be a continuous positive function with support contained in V such that \(\int\varphi_{V}\mu=1\) . Let B be a base of the filter of compact neighborhoods of e.

Let \(g\) be in \(G\) and \(f\) be in \(L^1(G)\) . We have \((\varphi_V * \varepsilon_g) * f \to \varepsilon_g * f\) in \(L^1(G)\) along the filter of sections of \(\mathfrak{B}\) (INT, VIII, p. 172, § 4, no. 7, prop. 20), hence \(\widetilde{\pi}(\varphi_V * \varepsilon_g)\widetilde{\pi}(f)\) converges to \(\widetilde{\pi}(\varepsilon_g * f)\) in \(\mathcal{L}(E)\) . This implies that \(\widetilde{\pi}(\varphi_V * \varepsilon_g)\) converges simply along the filter of sections of \(\mathfrak{B}\) to a linear map \(\pi(g)\) from \(F\) to \(F\) ; this map is continuous by EVT, III, p. 26, cor. 3, since \(\| \widetilde{\pi}(\varphi_V * \varepsilon_g) \| \leqslant 1\) for every compact neighborhood \(V\) of \(e\) . There therefore exists a unique endomorphism of \(E\) which induces \(\pi(g)\) by passage to the subspaces. We still denote by \(\pi(g)\) this endomorphism; we have \(\| \pi(g) \| \leqslant 1\) .

Let \(f \in \mathrm{L}^1(\mathrm{G})\) . Since \(\widetilde{\pi}(\varphi_{\mathrm{V}} * \varepsilon_g) \widetilde{\pi}(f)\) converges to \(\widetilde{\pi}(\varepsilon_g * f)\) , we have \(\pi(g) \widetilde{\pi}(f) = \widetilde{\pi}(\varepsilon_g * f)\) .

For \(g = e\) , this relation shows that \(\pi(e)\) is the identity on \(F\) , hence \(\pi(e) = 1_{E}\) . Let \(g_1\) and \(g_2\) be in \(G\) . We have

\[
\pi (g _ {1}) \pi (g _ {2}) \widetilde {\pi} (f) = \pi (g _ {1}) \widetilde {\pi} (\varepsilon_ {g _ {2}} * f) = \widetilde {\pi} (\varepsilon_ {g _ {1} g _ {2}} * f) = \pi (g _ {1} g _ {2}) \widetilde {\pi} (f),
\]

whence \(\pi(g_1)\pi(g_2) = \pi(g_1g_2)\) on F, and hence on E. This proves that the map \(g \mapsto \pi(g)\) is a representation of G in E. Since \(\|\pi(g)\| \leqslant 1\) and \(\|\pi(g)^{-1}\| = \|\pi(g^{-1})\| \leqslant 1\) , the endomorphisms \(\pi(g)\) of E are isometric, hence unitary (EVT, V, p. 40, prop. 3).

Let \(f \in \mathrm{L}^{1}(\mathrm{G})\) and \(x \in E\) . The map \(g \mapsto \varepsilon_{g} * f\) from G into \(\mathrm{L}^{1}(\mathrm{G})\) is continuous (INT, VIII, p. 136, § 2, no. 5 and p. 144, § 3, no. 2, formula (5)), and \(\widetilde{\pi}\) is continuous, hence the map \(g \mapsto \widetilde{\pi}(\varepsilon_{g} * f)x = \pi(g)\widetilde{\pi}(f)x\) is continuous at \(g = e\) . Since F is dense in E, the representation \(\pi\) is unitary (lemma 4 of V, p. 380).

Let \(f_1\) and \(f_2\) be in \(L^1(G)\) . According to INT, VIII, p. 127, §1, no. 5, prop. 7, we have the relation

\[
f _ {1} * f _ {2} = \int_ {\mathrm{G}} f _ {1} (g) (\varepsilon_ {g} * f _ {2}) d \mu (g)
\]

in L \(^{1}\) (G), whence

\[
\begin{array}{l} \widetilde {\pi} (f _ {1}) \widetilde {\pi} (f _ {2}) = \widetilde {\pi} (f _ {1} * f _ {2}) = \int_ {\mathrm{G}} f _ {1} (g) \widetilde {\pi} (\varepsilon_ {g} * f _ {2}) d \mu (g) \\ \qquad = \int_ {\mathrm{G}} f _ {1} (g) \pi (g) \widetilde {\pi} (f _ {2}) d \mu (g) = \Big (\int_ {\mathrm{G}} f _ {1} (g) \pi (g) d \mu (g) \Big) \widetilde {\pi} (f _ {2}), \end{array}
\]

using INT, VI, p. 9, § 1, no. 1, prop. 1. It follows that

\[
\widetilde {\pi} (f) = \int_ {\mathrm{G}} f (g) \pi (g) d \mu (g) = \pi (f)
\]

for all \(f \in \mathrm{L}^{1}(\mathrm{G})\) . The proposition is proved.

PROPOSITION 13. --- Let \(\varphi\in\mathrm{L}^{\infty}(\mathrm{G})\) . Denote by \(\lambda_{\varphi}\) the continuous linear form \(f\mapsto\langle f,\varphi\rangle\) on \(\mathrm{L}^{1}(\mathrm{G})\) . Then \(\varphi\) is the class of a function of positive type on G if and only if \(\lambda_{\varphi}\) is a continuous positive linear form on \(\mathrm{L}^{1}(\mathrm{G})\) .

Suppose that \(\lambda_{\varphi}\) is a continuous positive linear form on \(L^{1}(G)\) . Let \((E, x, \widetilde{\pi})\) be a cyclic Hilbertian realization of \(\lambda\) (prop. 5 of V, p. 438).

Let \(\pi\) be a unitary representation of G in E such that \(\pi(f) = \widetilde{\pi}(f)\) for all \(f \in \mathrm{L}^1(\mathrm{G})\) (prop. 12). We then have \(\lambda_{\varphi}(f) = \langle x | \pi(f)x \rangle\) for all \(f \in \mathrm{L}^1(\mathrm{G})\) .

Since

\[
\pi (f) = \int_ {\mathrm{G}} f (g) \pi (g) d \mu (g)
\]

for all \(f\in \mathrm{L}^1 (\mathrm{G})\) , it follows that

\[
\int_ {\mathrm{G}} f (g) \varphi (g) d \mu (g) = \lambda_ {\varphi} (f) = \langle x | \pi (f) x \rangle = \int_ {\mathrm{G}} f (g) \langle x | \pi (g) x \rangle d \mu (g)
\]

for all \(f \in \mathrm{L}^{1}(\mathrm{G})\) . Thus, \(\varphi\) is the class in \(\mathrm{L}^{\infty}(\mathrm{G})\) of the function on G defined by \(g \mapsto \langle x | \pi(g)x \rangle\) , which is a function of positive type on G (Prop. 9 of V, p. 443).

Conversely, let \(\varphi\in\mathrm{Pos}(\mathrm{G})\) . Let \(f\in\mathcal{K}(\mathrm{G})\) . We then have

\[
\begin{array}{c} \lambda_ {\varphi} (f ^ {*} * f) = \int_ {G} \varphi (y) \Big (\int_ {G} \Delta (z) ^ {- 1} \overline {{f (z ^ {- 1})}} f (z ^ {- 1} y) d \mu (z) \Big) d \mu (y) \\ = \int_ {G} \varphi (y) \Big (\int_ {G} \overline {{f (z)}} f (z y) d \mu (z) \Big) d \mu (y) \end{array}
\]

The continuous function on G × G defined by \((z, y) \mapsto \varphi(y)\overline{f(z)}f(zy)\) is bounded and

\[
\begin{array}{c} \int_ {\mathrm{G}} ^ {*} | \varphi (y) | \Big (\int_ {\mathrm{G}} ^ {*} | \overline {{f (z)}} f (z y) | d \mu (z) \Big) d \mu (y) \\ \leqslant \varphi (e) \int_ {\mathrm{G}} ^ {*} \int_ {\mathrm{G}} ^ {*} | f (z) | | f (z y) | d \mu (z) d \mu (y) \\ = \varphi (e) \Big (\int_ {\mathrm{G}} ^ {*} | f (z) | d \mu (z) \Big) ^ {2} \end{array}
\]

By INT, V, p. 94, § 8, no. 3, prop. 8, we therefore deduce from the Lebesgue–Fubini theorem (INT, V, p. 96, § 8, no. 4, th. 1) that

\[
\begin{array}{c} \lambda_ {\varphi} (f ^ {*} * f) = \int_ {\mathrm{G}} \overline {{f (z)}} \Bigl (\int_ {\mathrm{G}} \varphi (y) f (z y) d \mu (y) \Bigr) d \mu (z) \\ = \int_ {\mathrm{G} \times \mathrm{G}} \overline {{f (z)}} f (y) \varphi (z ^ {- 1} y) d (\mu \otimes \mu) (z, y). \end{array}
\]

This implies that \(\lambda_{\varphi}(f^{*} * f) \geqslant 0\) by Theorem 1, (i) of V, p. 432 applied to the measure induced by \(\mu\) on the support of \(f\) (INT, IV, p. 186, § 5, n° 7, def. 4). From this one deduces by continuity that \(\lambda_{\varphi}(f^{*} * f) \geqslant 0\) for all \(f \in \mathrm{L}^1(\mathrm{G})\) .

The spaces \(\mathrm{L}^{\infty}(\mathrm{G})\) and \(\mathcal{C}_{b}(\mathrm{G})\) are equipped with their Banach space topologies, and we denote the norm by \(f \mapsto \|f\|_{\infty}\) . We shall call here the weak topology on \(\mathrm{L}^{\infty}(\mathrm{G})\) , \(\mathcal{C}_{b}(\mathrm{G})\) or \(\operatorname{Pos}(\mathrm{G})\) the topology induced by the weak topology \(\sigma(\mathrm{L}^{\infty}(\mathrm{G}), \mathrm{L}^{1}(\mathrm{G}))\) .

Since \(L^{\infty}(G)\) is identified with the dual of \(L^{1}(G)\) (INT, V, p. 61, §5, no. 8, th. 4), every closed ball of \(L^{\infty}(G)\) is compact for the weak topology (EVT, III, p. 17, cor. 3).

COROLLARY. --- In \(\mathrm{L}^{\infty}(\mathrm{G})\) equipped with the weak topology, the set \(\operatorname{Pos}(\mathrm{G})\) is closed, and the set \(\operatorname{Pos}_{\leqslant 1}(\mathrm{G})\) is compact.

The first assertion follows from the proposition, and the second then follows from EVT, III, loc. cit.

In general, \(\mathrm{Pos}_{1}(\mathrm{G})\) is not compact for the weak topology, as shown by the example of the group R (exercise 10 of V, p. 497).

PROPOSITION 14. --- The set of extreme points of \(\mathrm{Pos}_{\leqslant 1}(\mathrm{G})\) is equal to the union of the set of extreme points of \(\mathrm{Pos}_1(\mathrm{G})\) and of the zero function. Moreover, the closed convex hull of the set of extreme points of \(\mathrm{Pos}_1(\mathrm{G})\) contains \(\mathrm{Pos}_1(\mathrm{G})\) .

By the corollary of Proposition 13, the set \(\mathrm{Pos}_{\leqslant 1}(\mathrm{G})\) is a cap of the pointed convex cone \(\mathrm{Pos}(\mathrm{G})\) , which is defined by the gauge \(\varphi \mapsto \varphi(e)\) (EVT, II, p. 61, def. 3 and prop. 4). Its extreme points are therefore the zero function and the elements \(\varphi\) of \(\mathrm{Pos}_1(\mathrm{G})\) belonging to the extreme generators of \(\mathrm{Pos}(\mathrm{G})\) (EVT, II, p. 62, cor. 1). These are the extreme points of \(\mathrm{Pos}_1(\mathrm{G})\) (EVT, II, p. 61, prop. 3). The first assertion follows from this.

Let us prove the second assertion. Let \(\varphi\in\mathrm{Pos}_{1}(\mathrm{G})\) . Since the set \(\mathrm{Pos}_{\leqslant 1}(\mathrm{G})\) is compact for the weak topology (corollary of Prop. 13), there exists a filter \(\mathfrak{F}\) on the convex hull of the extreme points of \(\mathrm{Pos}_{\leqslant 1}(\mathrm{G})\) which converges to \(\varphi\) (EVT, II, p. 59, th. 1). As the closed balls of the Banach space \(\mathrm{L}^{\infty}(\mathrm{G})\) are closed for the weak topology and since \(\|\varphi\|_{\infty}=1\) , we have \(\lim_{\psi,\mathfrak{F}}\|\psi\|_{\infty}=1\) (indeed, otherwise there would exist a real number \(c<1\) and a filter \(\mathfrak{G}\) on the closed ball of radius \(c\) in \(\mathrm{L}^{\infty}(\mathrm{G})\) which would be finer than \(\mathfrak{F}\) , which would imply that \(\varphi\) belongs to this closed ball).

According to the description of the extreme points of \(\mathrm{Pos}_{\leqslant1}(\mathrm{G})\) , every element \(\psi\) of the convex hull of the extreme points of \(\mathrm{Pos}_{\leqslant1}(\mathrm{G})\) is a function of positive type on G of the form

\[
\psi = \sum_ {i \in \mathrm{I}} t _ {i} \psi_ {i},
\]

where I is a finite set, \(\psi_{i}\) is an extreme point of \(\mathrm{Pos}_{1}(\mathrm{G})\) for all \(i \in I\) , and \(t_{i} \in [0,1]\) . We have \(\sum_{i} t_{i} = \psi(e) = \|\psi\|_{\infty} \leqslant 1\) (lemma 4 of V, p. 446). If \(\psi \neq 0\) , the function

\[
\frac {\psi}{\| \psi \| _ {\infty}} = \sum_ {i \in \mathrm{I}} \frac {t _ {i}}{\psi (e)} \psi_ {i}
\]

therefore belongs to the convex hull of the extreme points of \(\mathrm{Pos}_{1}(\mathrm{G})\) . Since \(\lim_{\psi,\mathfrak{F}}\|\psi\|_{\infty}^{-1}\psi=\varphi\) , we conclude that \(\varphi\) belongs to the closed convex hull of the extreme points of \(\mathrm{Pos}_{1}(\mathrm{G})\) . The assertion follows.

